\section{Writing back}\label{sec:writeback}

A memory that agents only read is a library. What makes the brain a collaborator is that agents write to it, and the dangerous part of the design is there: a store that a hundred agents can write to will fill with confident errors unless the path in is narrow. We use three gates (a packet gate on recall, a verification router on truth, and quarantine for anything generated by the system itself) and one rule, that nothing is ever deleted.

\subsection{Packets}
An agent that has a result writes a \emph{packet}: a small bundle (Definition~\ref{def:bundle}) of new cards and links, each card with its statement, its exact hypotheses, its proof text or certificate, its source, and the links that place it among existing cards. A packet is typically a handful of cards: a lemma, the counterexample that shows a hypothesis cannot be dropped, the refuted route that motivated both, and the links among them and to the cards they use.

Before a packet lands it passes a \emph{recall gate}. A held-out set of recall questions, disjoint from every question used to tune the ranking, is run against a copy of the store with the packet added and against the same copy without it. The packet lands only if it costs no answer that the store without it gave. The gate exists because a new card can displace an old one from the first eight results of a question it does not answer; with thirteen thousand cards and hundreds of packets a week, small displacements add up. The gate measures recall only. It says nothing about whether the packet is true, which is the next gate's business.

\subsection{Truth status}
Every card enters as \textsf{unverified}. A status changes only through the verification router, which decomposes a claim into obligations and asks for evidence suited to each: arithmetic reproduction; symbolic identity under explicit domains; randomized or exhaustive finite testing; human proof review; formal proof checking; review of the fidelity of a formalization to the informal statement; citation entailment with locators; and independent reproduction. No single kind of evidence is accepted as universal; promotion requires coverage at the claim's actual scope.

Evidence rows alone are never authority. A status change requires a receipt: an HMAC-SHA-256 tag, under a key held outside the store and fixed to one role, over the exact digests of the evidence, the claim revision, the artifact, the verifier and its toolchain, the source, and a single-use nonce. Human review of fidelity, automated checking and independent reproduction use distinct roles, and a key configured for one role cannot sign another role's evidence. Evidence, attestation and receipt are appended in one transaction. The status of a card is not a stored field that anyone can write: it is recomputed by replaying every verification event and re-authenticating its evidence, so a status written directly into the database is detected as an error.

The Lean path is concrete. The Lean authority snapshots a dependency-free local module graph, runs only the hash-pinned Lean executable of the installed toolchain on every module, requires a report that no named declaration uses an axiom beyond the standard ones, and signs only a fresh seal bound to the claim digest. Reports produced by agents themselves (``this was checked in Lean'') are kept as discovery reports and can never be promoted to evidence; a trusted local verifier must rerun the artifact.

In practice most cards remain \textsf{unverified} or \textsf{provisional}, and the trust line printed under each card (Section~\ref{sec:hook}) says so. This is the right state of affairs for a research notebook: the store records what was claimed, by what route, with what support, and what later refuted it.

\subsection{Mistakes}
A refuted route is as useful to the next agent as a proof, and often more. A mistake record stores the minimal failing step, its cause, the hypotheses it affects, the counterexample or repair, and a prevention action phrased as a concrete check (``division: prove the denominator nonzero''). It is joined to the claims it concerns by negative links (\texttt{fails\_under}), so that spreading and walking show it as a caution beside the claim. A reflex record pairs explicit activation cues with ordered actions, required evidence and a stopping condition. Neither changes the status of any claim. The store holds $160$ \texttt{fails\_under} and $138$ \texttt{refutes} links; $114$ of the former were added later than both of their endpoints, that is, an agent found a pitfall of an existing result and recorded it.

\subsection{Nothing is deleted}
A wrong card is superseded, not deleted; a wrong link is retracted, and the retraction is kept with its reason. On 22 October the store held $1{,}242$ \texttt{supersedes} links and $231$ retracted links. The rule costs disk and occasionally clutters a printout with an old revision, and in exchange the history of every result is complete: an agent can see that a statement was once claimed without a hypothesis, who found the counterexample, and what the corrected statement is.

\subsection{Consolidation}
Between campaigns, a consolidation pass replays the store and proposes work: contradictions to check, claims with missing evidence, cards with no links, hubs to audit, mistakes to rehearse, and analogies suggested by shared formula structure. Every proposal is an object of a type whose status can only be \textsf{quarantined} and whose field \texttt{can\_promote\_truth} is false; treating one as evidence raises an error. A proposal becomes knowledge only when an agent does the proposed work and lands the result through the gates above. The same contamination guard used in evaluation (Section~\ref{sec:tests}) is applied before consolidation, so that a test answer cannot re-enter the store through a generated association.
